%%% Local Variables:
%%% mode: latex
%%% TeX-master: t
%%% End:

% ============================================================
% 第三章：数学公式与定理
% 本章介绍 LaTeX 数学排版，包括公式、符号速查和定理环境。
% ============================================================

\chapter{数学公式与定理}\label{chap:math}

\section{行内公式}

行内公式嵌入在正文中，用一对美元符号 \verb|$ $| 包围：

\begin{verbatim}
质子的静止质量为 $m_p = 938.3$ MeV/$c^2$。
\end{verbatim}

效果：质子的静止质量为 $m_p = 938.3$ MeV/$c^2$。

常用行内公式写法：
\begin{compactitem}
    \item 下标：\verb|$a_i$| → $a_i$，上标：\verb|$E^2$| → $E^2$
    \item 分数（行内）：\verb|$m/c^2$| → $m/c^2$（行内分数建议用斜线形式）
    \item 希腊字母：\verb|$\alpha, \beta, \gamma$| → $\alpha, \beta, \gamma$
    \item 向量（粗体）：\verb|$\mathbf{p}$| → $\mathbf{p}$
\end{compactitem}

\section{行间公式}

行间公式独立成行，用 \texttt{equation} 环境，自动按章节编号：

\begin{verbatim}
\begin{equation}
    E = mc^2
    \label{eq:einstein}
\end{equation}
\end{verbatim}

\begin{equation}
    E = mc^2
    \label{eq:einstein}
\end{equation}

由式~(\ref{eq:einstein})~可知，质能等价。

更复杂的示例——薛定谔方程：
\begin{equation}
    i\hbar \frac{\partial}{\partial t} \Psi(\mathbf{r}, t)
    = \left[ -\frac{\hbar^2}{2m} \nabla^2 + V(\mathbf{r}, t) \right] \Psi(\mathbf{r}, t)
    \label{eq:schrodinger}
\end{equation}

\textbf{格式规范}：按格式要求，公式序号右端对齐，
\texttt{equation} 环境已自动处理。
公式中第一次出现的物理量符号应在公式后用"式中"段落注释。

\subsection{公式后的符号注释}

格式规范要求对公式中的物理量进行说明，写法如下：

\begin{verbatim}
\begin{equation}
    \sigma = \frac{I_0 - I}{I_0 N t}
    \label{eq:cross_section}
\end{equation}

式中：$\sigma$——反应截面（cm$^2$）；\\
\phantom{式中：}$I_0$——入射粒子数；\\
\phantom{式中：}$I$——出射粒子数；\\
\phantom{式中：}$N$——靶核数密度（cm$^{-3}$）；\\
\phantom{式中：}$t$——靶厚（cm）。
\end{verbatim}

\begin{equation}
    \sigma = \frac{I_0 - I}{I_0 N t}
    \label{eq:cross_section}
\end{equation}

式中：$\sigma$——反应截面（cm$^2$）；\\
\phantom{式中：}$I_0$——入射粒子数；\\
\phantom{式中：}$I$——出射粒子数。

\section{多行公式}

当公式较长需要换行，或需要对齐多个等式时，使用 \texttt{align} 环境：

\begin{verbatim}
\begin{align}
    Q &= (M_A + M_a - M_B - M_b)c^2  \label{eq:q_value} \\
      &= T_B + T_b - T_a              \label{eq:q_kinetic}
\end{align}
\end{verbatim}

\begin{align}
    Q &= (M_A + M_a - M_B - M_b)c^2  \label{eq:q_value} \\
      &= T_B + T_b - T_a              \label{eq:q_kinetic}
\end{align}

\verb|&| 表示对齐位置，\verb|\\| 换行。每行都有独立编号。

如果不需要某行编号，在该行末尾加 \verb|\nonumber|：
\begin{verbatim}
\begin{align}
    f(x) &= x^2 + 2x + 1  \nonumber \\
         &= (x+1)^2       \label{eq:factor}
\end{align}
\end{verbatim}

\section{矩阵}

\begin{verbatim}
\begin{equation}
    \mathbf{A} = \begin{pmatrix}
        a_{11} & a_{12} & \cdots & a_{1n} \\
        a_{21} & a_{22} & \cdots & a_{2n} \\
        \vdots & \vdots & \ddots & \vdots \\
        a_{m1} & a_{m2} & \cdots & a_{mn}
    \end{pmatrix}
\end{equation}
\end{verbatim}

\begin{equation}
    \mathbf{A} = \begin{pmatrix}
        a_{11} & a_{12} & \cdots & a_{1n} \\
        a_{21} & a_{22} & \cdots & a_{2n} \\
        \vdots & \vdots & \ddots & \vdots \\
        a_{m1} & a_{m2} & \cdots & a_{mn}
    \end{pmatrix}
\end{equation}

矩阵括号类型：
\begin{compactitem}
    \item \texttt{pmatrix}：圆括号 $(\,)$
    \item \texttt{bmatrix}：方括号 $[\,]$
    \item \texttt{vmatrix}：竖线（行列式）$|\,|$
    \item \texttt{Bmatrix}：花括号 $\{\,\}$
\end{compactitem}

\section{常用数学符号速查}

\subsection{希腊字母}

\begin{table}[H]
\centering
\bicaption{常用希腊字母}{Common Greek letters}
\begin{tabular}{llll|llll}
\hline
命令 & 符号 & 命令 & 符号 & 命令 & 符号 & 命令 & 符号 \\
\hline
\verb|\alpha|   & $\alpha$   & \verb|\Alpha|   & $A$       & \verb|\beta|    & $\beta$    & \verb|\gamma|   & $\gamma$ \\
\verb|\delta|   & $\delta$   & \verb|\Delta|   & $\Delta$  & \verb|\epsilon| & $\epsilon$ & \verb|\theta|   & $\theta$ \\
\verb|\lambda|  & $\lambda$  & \verb|\Lambda|  & $\Lambda$ & \verb|\mu|      & $\mu$      & \verb|\nu|      & $\nu$ \\
\verb|\pi|      & $\pi$      & \verb|\Pi|      & $\Pi$     & \verb|\sigma|   & $\sigma$   & \verb|\Sigma|   & $\Sigma$ \\
\verb|\phi|     & $\phi$     & \verb|\Phi|     & $\Phi$    & \verb|\psi|     & $\psi$     & \verb|\Psi|     & $\Psi$ \\
\verb|\omega|   & $\omega$   & \verb|\Omega|   & $\Omega$  & \verb|\eta|     & $\eta$     & \verb|\xi|      & $\xi$ \\
\hline
\end{tabular}
\end{table}

\subsection{常用运算符}

\begin{compactitem}
    \item 积分：\verb|\int_0^\infty f(x)dx| → $\int_0^\infty f(x)dx$
    \item 求和：\verb|\sum_{i=1}^{n} a_i| → $\sum_{i=1}^{n} a_i$
    \item 乘积：\verb|\prod_{i=1}^{n}| → $\prod_{i=1}^{n}$
    \item 分数：\verb|\frac{分子}{分母}| → $\frac{a+b}{c}$
    \item 根号：\verb|\sqrt{x}|, \verb|\sqrt[3]{x}| → $\sqrt{x}$, $\sqrt[3]{x}$
    \item 偏微分：\verb|\partial| → $\partial$
    \item 约等于：\verb|\approx| → $\approx$；正比于：\verb|\propto| → $\propto$
    \item 向量算符：\verb|\nabla| → $\nabla$；\verb|\hbar| → $\hbar$
\end{compactitem}

\section{定理、引理与证明环境}

模板在 \texttt{ustcextra.sty} 中预定义了以下环境：

\begin{verbatim}
定理:     \begin{theorem}...\end{theorem}
引理:     \begin{lemma}...\end{lemma}
推论:     \begin{corollary}...\end{corollary}
定义:     \begin{definition}...\end{definition}
命题:     \begin{proposition}...\end{proposition}
注:       \begin{remark}...\end{remark}
证明:     \begin{proof}...\end{proof}
\end{verbatim}

示例——一个简单的定理及其证明：

\begin{theorem}
    对于任意实数 $a$ 和 $b$，有 $(a+b)^2 = a^2 + 2ab + b^2$。
\end{theorem}

\begin{proof}
    展开左边：$(a+b)^2 = (a+b)(a+b) = a^2 + ab + ba + b^2 = a^2 + 2ab + b^2$，
    故等式成立。
\end{proof}

定理类环境按章节自动编号。在 \texttt{main.tex} 中可以用 \verb|\newtheorem| 自定义新的定理类环境。

\section{本章小结}

本章介绍了行内公式、行间公式、多行公式对齐、矩阵写法、
常用数学符号以及定理环境的使用方法。
第~\ref{chap:floats}~章将介绍图片和表格的插入与排版。
